% Candidate teaching additions drawn from 04-lec.pdf.
% This block appears before the existing Section 3.1 slides so it can be
% reviewed as a unit. It may later be distributed among Sections 3.1 and 3.2.
\section[Review additions]{Additional Details for Review}
\sectioncard{Review}{Additional Details}

\begin{frame}{Forming a Minor}
To form $A_{32}$, delete row $3$ and column $2$.
\[
A=\begin{bmatrix}
1&\textcolor{lectureblue}{2}&5&0\\
2&\textcolor{lectureblue}{0}&4&1\\
\textcolor{red}{3}&\textcolor{red}{1}&\textcolor{red}{0}&\textcolor{red}{7}\\
0&\textcolor{lectureblue}{4}&2&0
\end{bmatrix}
\qquad\Longrightarrow\qquad
A_{32}=\begin{bmatrix}
1&5&0\\
2&4&1\\
0&2&0
\end{bmatrix}.
\]
\begin{center}
\textcolor{red}{Red entries belong to row $3$.}
\qquad
\textcolor{lectureblue}{Blue entries belong to column $2$.}
\end{center}
\pause
Everything that remains keeps its original order.
\end{frame}

\begin{frame}{From a Minor to a Cofactor}
Continue with the same position, $(3,2)$.
\begin{enumerate}
  \item The minor is
  \[
  M_{32}=\det A_{32}
  =\begin{vmatrix}1&5&0\\2&4&1\\0&2&0\end{vmatrix}=-2.
  \]
  \pause
  \item The checkerboard sign at $(3,2)$ is negative:
  \[
  C_{32}=(-1)^{3+2}M_{32}=-(-2)=2.
  \]
  \pause
  \item Since $a_{32}=1$, this entry contributes
  \[
  a_{32}C_{32}=1(2)=2
  \]
  to an expansion across row $3$ or down column $2$.
\end{enumerate}
\end{frame}

\begin{frame}{Choosing a Computation Method}
\small
\renewcommand{\arraystretch}{1.35}
\begin{center}
\begin{tabular}{>{\raggedright\arraybackslash}p{.40\linewidth}>{\raggedright\arraybackslash}p{.48\linewidth}}
\toprule
Matrix structure & Efficient method \\
\midrule
$2\times2$ & Use $ad-bc$ \\
One row or column has many zeros & Use cofactor expansion there \\
Triangular & Multiply the diagonal entries \\
Large and dense & Use elimination and track determinant changes \\
\bottomrule
\end{tabular}
\end{center}
\pause
A direct cofactor calculation grows roughly like $n!$, while elimination takes about
\[
\frac{2n^3}{3}
\]
arithmetic operations. For $n=25$, these scales are approximately $1.55\times10^{25}$ and $10{,}417$.
\end{frame}

\begin{frame}{Operation Ledger: Reduction}
\small
Compute the determinant of
\[
A=\begin{bmatrix}
0&2&1\\
1&1&0\\
2&5&3
\end{bmatrix}.
\]
Use the following operations:
\[
\begin{aligned}
R_1&\leftrightarrow R_2,\\
R_3&\leftarrow R_3-2R_1,\\
R_2&\leftarrow \tfrac12 R_2,\\
R_3&\leftarrow R_3-3R_2.
\end{aligned}
\]
\pause
The resulting triangular matrix is
\[
U=\begin{bmatrix}
1&1&0\\
0&1&\tfrac12\\
0&0&\tfrac32
\end{bmatrix}.
\]
\end{frame}

\begin{frame}{Operation Ledger: Recovering $\det A$}
\small
\renewcommand{\arraystretch}{1.35}
\begin{center}
\begin{tabular}{lc}
\toprule
Operation & Effect on the current determinant \\
\midrule
One row interchange & multiply by $-1$ \\
Two row replacements & multiply by $1$ \\
$R_2\leftarrow\tfrac12R_2$ & multiply by $\tfrac12$ \\
\bottomrule
\end{tabular}
\end{center}
Therefore
\[
\det U=(-1)\left(\frac12\right)\det A.
\]
\pause
Since
\[
\det U=1(1)\left(\frac32\right)=\frac32,
\]
we obtain
\[
\boxed{\det A=-2\det U=-3}.
\]
\end{frame}

\begin{frame}{Combining Elimination and Cofactor Expansion}
Let
\[
A=\begin{bmatrix}
0&1&2&-1\\
2&5&7&3\\
0&3&6&2\\
2&5&4&-2
\end{bmatrix}.
\]
Apply the row replacement $R_4\leftarrow R_4-R_2$:
\[
\det A=
\begin{vmatrix}
0&1&2&-1\\
2&5&7&3\\
0&3&6&2\\
0&0&-3&-5
\end{vmatrix}.
\]
\pause
Now the first column contains a single nonzero entry.
\end{frame}

\begin{frame}{Combining Methods: Finish}
Expand down the first column:
\[
\det A
=2(-1)^{2+1}
\begin{vmatrix}
1&2&-1\\
3&6&2\\
0&-3&-5
\end{vmatrix}.
\]
\pause
The $3\times3$ determinant is
\begin{align*}
&1[-30+6]-2[-15-0]+(-1)[-9-0]\\
&=-24+30+9=15.
\end{align*}
Therefore
\[
\boxed{\det A=-2(15)=-30}.
\]
The row replacement created the zero pattern, and the cofactor expansion used it.
\end{frame}

\begin{frame}{Linearity in One Row}
Hold every row except one fixed. In that one row, the determinant satisfies both linearity rules:
\[
D(\ldots,c\vect u,\ldots)=cD(\ldots,\vect u,\ldots),
\]
\[
D(\ldots,\vect u+\vect v,\ldots)
=D(\ldots,\vect u,\ldots)+D(\ldots,\vect v,\ldots).
\]
\pause
The same rules hold when one column varies and all other columns remain fixed.
\medskip

\begin{alertblock}{Keep the other rows or columns fixed}
These rules do not imply $\det(A+B)=\det A+\det B$ when every row changes at once.
\end{alertblock}
\end{frame}

\begin{frame}{Concept Check}
Let $A$ be a $3\times3$ matrix with $\det A=5$.
\begin{enumerate}
  \item One row of $A$ is multiplied by $2$. What is the new determinant?
  \item The entire matrix is multiplied by $2$. What is $\det(2A)$?
  \item Three row interchanges are performed. What is the new determinant?
\end{enumerate}
\pause
\[
\boxed{10},
\qquad
\boxed{2^3(5)=40},
\qquad
\boxed{(-1)^3(5)=-5}.
\]
\pause
Scaling one row and scaling the whole matrix are different operations.
\end{frame}
